\documentclass[12pt]{article}
\usepackage{pmmeta}
\pmcanonicalname{LeanCatBenchmark}
\pmcreated{2026-03-01 17:50:00}
\pmmodified{2026-03-01 17:50:00}
\pmowner{codex}{0}
\pmmodifier{codex}{0}
\pmtitle{LeanCat: a benchmark for formal category theory}
\pmrecord{0}{0}
\pmprivacy{1}
\pmauthor{codex}{0}
\pmtype{Article}
\pmclassification{msc}{18A99}
\pmclassification{msc}{68V20}
\pmrelated{ProofAssistant}
\pmrelated{LeanTheoremProver}
\pmrelated{AdjointFunctor}
\pmdefines{LeanCat benchmark}

\endmetadata

\usepackage{amsmath,amssymb,amsfonts}

\begin{document}
\section*{Overview}
\textsc{LeanCat} is a curated suite of formalization tasks designed to stress-test Lean's category-theory library (\texttt{mathlib}) together with language-model-guided automation.  Each task is a theorem statement written in Lean together with hints on the expected proof strategy.  The benchmark currently stores 208 tasks grouped into clusters: limits and colimits, adjunctions, monoidal structures, enriched categories, and higher categorical constructions.

\section*{Dataset geometry}
The tasks form a directed acyclic graph whose nodes are theorems and whose edges encode dependencies.  Each node carries:
\begin{itemize}
  \item the Lean level (a proxy for proof difficulty),
  \item tags for the categorical notions involved, and
  \item links to human-written reference proofs (when available).
\end{itemize}
The dependency graph is used to design curricula for autoformalization agents: a learner must solve prerequisite tasks before attempting harder ones.

\section*{Automation baselines}
The benchmark evaluates several strategies:
\begin{description}
  \item[LLM-guided tactics] combine large language models with Lean's tactic framework to synthesize proof scripts.
  \item[Retrieval-augmented proving] retrieves structurally similar proofs and adapts them to the goal via rewriting.
  \item[Curriculum learning] interleaves automated proof search with human hints along the dependency DAG.
\end{description}
The reported experiments show that purely neural tactics can solve easy limit lemmas, while Lean's native search still outperforms them on coherence-heavy statements such as triangle identities for adjunctions.

\section*{Outlook}
LeanCat highlights concrete obstacles for fully automated category-theoretic reasoning: missing typeclass instances, coherence lemmas for monoidal structures, and proof search in higher inductive contexts.  The benchmark therefore doubles as a roadmap for strengthening mathlib and for evaluating future autoformalization systems.

\end{document}
